\section{Ba lớp mô hình và vai trò của từng lớp}

Model thuật toán mô tả tuần tự bước OPS, SAC và SBM, giữ các
checkpoint để có thể tìm sai khác theo vòng. Nó hữu ích khi
gỡ lỗi nhưng dễ lặp lại lỗi của RTL nếu cả hai dùng cùng
recurrence. Golden của bộ nhân chọn trực tiếp các vị trí
bit cao theo OPS và n, rồi nhân hai hệ số nguyên đã xác định.
Decoder của golden làm việc với mã Posit thô. Cách này kiểm
kết quả của RTL mà không cần giả định pipeline hoặc cách
token dịch sang T2.

Ideal dùng tất cả bit fraction và tính bằng số nguyên.
Với một giao dịch độc lập, ideal là tổng $AB+C$ trước RNE.
Với chuỗi, ideal là tổng các tích đủ fraction theo cả dãy,
không thay state bằng một mã Posit sau mỗi bước. Golden và
ideal vì vậy trả lời hai câu hỏi khác nhau. Golden dùng
để phân loại một sai lệch là lỗi thực hiện hay đúng theo
đặc tả. Ideal là mốc để lượng hóa sai số của chính đặc tả
xấp xỉ và lượng tử hóa.

\begin{engineeringtable}{Mốc so sánh theo ranh giới quan sát.}{tab:oracle-boundaries}
\begin{tabularx}{\linewidth}{@{}>{\raggedright\arraybackslash}p{32mm}Y Y@{}}
\toprule
Điểm quan sát & Golden hoặc trạng thái mô hình & Đại lượng sai số \\
\midrule
Core trước Packer & Tích hệ số chọn bit, $M$, dấu, scale, số vòng và cờ cắt. & Hai mốc $P_{12}$ và $P$ đầy đủ. \\
NF retire & $R(R(\widehat P)+C)$ với mã đã nhận. & So giá trị decode của D với $AB+C$. \\
F retire & $R(\widehat P+C)$ theo hợp đồng FULL\_TERM. & So với cùng tổng $AB+C$ ở phép đo paired. \\
Quire ack/READ & Trạng thái Q bigint, count, cờ, mã ở READ. & So kết quả cuối với tổng tích đầy đủ của cả chuỗi. \\
\bottomrule
\end{tabularx}
\end{engineeringtable}

Golden cộng trên lưới rộng trước encode, để tránh mất phần
dư ở các cặp có chênh lệch scale lớn. Đối với NF, golden
đóng gói tích trước rồi giải mã lại mã tích để cộng. Đối với
F, cùng hệ số Q24 được đưa thẳng vào tổng chính xác theo hợp
đồng. Các phép kiểm trạng thái quire dùng bigint độc lập
trên lưới Q và kiểm sau mỗi response. Điểm pack để đo sai
số chuỗi là READ cuối.

\section{Driver, scoreboard và monitor}

\blockfigure{Luồng kiểm chứng. Scoreboard chốt expected tại accept và so đúng một lần tại retire.}{fig:scoreboard}{
  \node[block] (fixture) {Corpus hoặc\\fixture};
  \node[block,right=5mm of fixture] (drv) {Driver\\valid/ready};
  \node[block,right=5mm of drv] (dut) {DUT\\NF/F/quire};
  \node[block,right=5mm of dut] (mon) {Monitor\\retire};
  \node[store,below=7mm of dut] (sb) {Scoreboard\\d và flags};
  \node[block,left=5mm of sb] (oracle) {Oracle\\số nguyên};
  \draw[data] (fixture)--(drv);
  \draw[data] (drv)--(dut);
  \draw[data] (dut)--(mon);
  \draw[data] (mon.south)|-(sb.east);
  \draw[control] (drv)--(oracle);
  \draw[data] (oracle)--(sb);
}

Trong hồi quy có feedback, scoreboard giữ state cho từng
lõi. Tại accept, nó tính expected từ A/B, C hiệu dụng,
OPS và n. Khi driver chọn hồi tiếp, C là kết quả đã retire
của chính lõi đang kiểm. Khi chạy paired standalone, mỗi lõi
nhận cùng C cố định từ corpus. Tách hai chế độ này tránh
so sánh sai khi hai chuỗi đã đi tới hai state khác nhau.

Monitor giao thức quan sát valid/ready ở cạnh clock. Khi
out valid mà ready thấp, kết quả vẫn thuộc về DUT và phải
giữ nguyên đến retire. Khi reset hủy request, scoreboard
bỏ expected thuộc epoch cũ. Sau khi reset nhả, một cửa sổ
quan sát không được có response cũ. Cuối lượt cần kiểm số
đếm accepted bằng completed cộng reset-aborted; đếm cả hai
vào completed sẽ làm sai khối lượng của ma trận.

\begin{figure}[!htbp]
\centering
\begin{tikzpicture}[x=.61cm,y=.65cm,>=Stealth,font=\scriptsize]
\draw[gray!28] (0,-.25)--(0,7.1) node[above,black] {0};
\draw[gray!28] (1,-.25)--(1,7.1) node[above,black] {1};
\draw[gray!28] (2,-.25)--(2,7.1) node[above,black] {2};
\draw[gray!28] (3,-.25)--(3,7.1) node[above,black] {3};
\draw[gray!28] (4,-.25)--(4,7.1) node[above,black] {4};
\draw[gray!28] (5,-.25)--(5,7.1) node[above,black] {5};
\draw[gray!28] (6,-.25)--(6,7.1) node[above,black] {6};
\draw[gray!28] (7,-.25)--(7,7.1) node[above,black] {7};
\draw[gray!28] (8,-.25)--(8,7.1) node[above,black] {8};
\draw[gray!28] (9,-.25)--(9,7.1) node[above,black] {9};
\draw[gray!28] (10,-.25)--(10,7.1) node[above,black] {10};
\draw[gray!28] (11,-.25)--(11,7.1) node[above,black] {11};
\draw[gray!28] (12,-.25)--(12,7.1) node[above,black] {12};
\draw[gray!28] (13,-.25)--(13,7.1) node[above,black] {13};
\draw[gray!28] (14,-.25)--(14,7.1) node[above,black] {14};
\draw[gray!28] (15,-.25)--(15,7.1) node[above,black] {15};
\draw[gray!28] (16,-.25)--(16,7.1) node[above,black] {16};
\draw[gray!28] (17,-.25)--(17,7.1) node[above,black] {17};
\node[anchor=east] at (-.3,7.3) {Cạnh};
\node[anchor=east] at (-.25,6.14) {accept};
\draw[blue!70!black,thick] (0,6.42)--(1,6.42)--(1,6)--(2,6)--(3,6)--(4,6)--(5,6)--(6,6)--(7,6)--(8,6)--(9,6)--(10,6)--(11,6)--(12,6)--(13,6)--(14,6)--(15,6)--(16,6)--(17,6)--(18,6);
\node[anchor=east] at (-.25,5.14) {pair\_retire};
\draw[blue!70!black,thick] (0,5)--(1,5)--(2,5)--(2,5.42)--(3,5.42)--(3,5)--(4,5)--(5,5)--(6,5)--(7,5)--(8,5)--(9,5)--(10,5)--(11,5)--(12,5)--(13,5)--(14,5)--(15,5)--(16,5)--(17,5)--(18,5);
\node[anchor=east] at (-.25,4.14) {T2 valid};
\draw[blue!70!black,thick] (0,4)--(1,4)--(2,4)--(3,4)--(4,4)--(5,4)--(5,4.42)--(6,4.42)--(7,4.42)--(7,4)--(8,4)--(9,4)--(10,4)--(11,4)--(12,4)--(13,4)--(14,4)--(15,4)--(16,4)--(17,4)--(18,4);
\node[anchor=east] at (-.25,3.14) {core\_done};
\draw[blue!70!black,thick] (0,3)--(1,3)--(2,3)--(3,3)--(4,3)--(5,3)--(6,3)--(7,3)--(7,3.42)--(8,3.42)--(8,3)--(9,3)--(10,3)--(11,3)--(12,3)--(13,3)--(14,3)--(15,3)--(16,3)--(17,3)--(18,3);
\node[anchor=east] at (-.25,2.14) {add\_valid};
\draw[blue!70!black,thick] (0,2)--(1,2)--(2,2)--(3,2)--(4,2)--(5,2)--(6,2)--(7,2)--(8,2)--(9,2)--(10,2)--(11,2)--(12,2)--(12,2.42)--(13,2.42)--(13,2)--(14,2)--(15,2)--(16,2)--(17,2)--(18,2);
\node[anchor=east] at (-.25,1.1400000000000001) {out\_valid};
\draw[blue!70!black,thick] (0,1)--(1,1)--(2,1)--(3,1)--(4,1)--(5,1)--(6,1)--(7,1)--(8,1)--(9,1)--(10,1)--(11,1)--(12,1)--(13,1)--(14,1)--(14,1.42)--(15,1.42)--(16,1.42)--(17,1.42)--(18,1.42);
\node[anchor=east] at (-.25,0.14) {out\_ready};
\draw[blue!70!black,thick] (0,0)--(1,0)--(2,0)--(3,0)--(4,0)--(5,0)--(6,0)--(7,0)--(8,0)--(9,0)--(10,0)--(11,0)--(12,0)--(13,0)--(14,0)--(15,0)--(16,0)--(17,0)--(17,0.42)--(18,0.42);
\draw[|<->|,teal!70!black,thin] (0,7.92)--(2,7.92) node[midway,above] {Parser};
\draw[|<->|,teal!70!black,thin] (2,7.92)--(7,7.92) node[midway,above] {Nhân (t=2)};
\draw[|<->|,teal!70!black,thin] (7,7.92)--(12,7.92) node[midway,above] {Adder S1--S5};
\draw[|<->|,teal!70!black,thin] (12,7.92)--(14,7.92) node[midway,above] {Pack};
\draw[|<->|,orange!80!black,thin] (14,7.92)--(17,7.92) node[midway,above] {Giữ (stall)};
\node[align=center] at (14.8,-.64) {Giữ d=1ae34800, flags=00 ở cạnh 14--16};
\end{tikzpicture}
\caption{Giản đồ xung trích trực tiếp từ trace mô phỏng RTL trên \texttt{ModelSim ASE 10.1d} cho ca thử \texttt{FT-05} của \texttt{MAC-F}. Mỗi cột là trạng thái trước cập nhật tại cạnh lên tương ứng. Với $t=2$, \texttt{out\_valid} lên ở cạnh $14 = 12+\max(1,2)$ (Parser $2$ + Nhân $5$ + Adder $5$ + Pack $2$), được giữ ổn định 3 chu kỳ khi \texttt{out\_ready=0} và hoàn tất \textit{retire} tại cạnh 17.}
\label{fig:ft05-wave}
\end{figure}


\begin{engineeringtable}{Trace ModelSim trên lõi \texttt{mac\_f\_wrapper}. Mốc cạnh tính từ cạnh \texttt{accept} của giao dịch đang xét.}{tab:control-trace}
\begin{tabularx}{\linewidth}{@{}>{\raggedright\arraybackslash}p{26mm}>{\raggedright\arraybackslash}p{28mm}Y@{}}
\toprule
Ca thử & Cạnh clock & Quan sát trạng thái \\
\midrule
FT-05, $n=12$ & 0, 2 & Nhận input, sau đó ba Parser retire vào core. \\
FT-05, $n=12$ & 7, 12 & \texttt{core\_done}, sau đó \texttt{add\_valid}. \\
Giữ kết quả & 14--16 & \texttt{out\_valid} cao, \texttt{out\_ready} thấp, \texttt{d=1ae34800}, flags \texttt{00}. \\
Retire & 17 & \texttt{out\_valid} và \texttt{out\_ready} cùng cao, đối chiếu payload. \\
FT-05, $n=0$ & 6, 14 & Core phát cơ số $M_0$, sau đó retire \texttt{1a410000}/\texttt{03}. \\
Reset hủy & 4; sau nhả & Reset assert khi đang xử lý; không có response cũ trong 30 cạnh quan sát. \\
Sau reset & 15 & Giao dịch kế tiếp retire chính xác bộ mẫu FT-05. \\
\bottomrule
\end{tabularx}
\end{engineeringtable}

Trace được lấy mẫu ngay trước cạnh lên đồng bộ của thanh ghi.
Cạnh \texttt{t2\_valid} mô tả token chuẩn bị được cộng dồn; giá trị accumulator
đã cập nhật xuất hiện ở chu kỳ kế tiếp. Bảng~\ref{tab:ft05-terms}
sử dụng thống nhất quy ước này để tránh gán nhầm $M$ cho token hiện tại.
Phiên mô phỏng đối chiếu giao thức nhận 4 giao dịch, retire 3 và hủy 1 do
\textit{reset}, khớp 100\% theo từng chu kỳ giữa hai phiên bản trước và sau tối ưu
tầng S1.

\section{Ma trận yêu cầu chức năng và kết quả xác nhận}

\begin{engineeringtable}{Tổng hợp đối chiếu yêu cầu thiết kế, bài kiểm thử và kết quả đạt được.}{tab:verification-matrix}
\begin{tabularx}{\linewidth}{@{}>{\raggedright\arraybackslash}p{28mm}Y Y@{}}
\toprule
Yêu cầu & Kịch bản kiểm thử và chỉ tiêu & Kết quả kiểm chứng \\
\midrule
Tích nội bộ Q24 & Bộ mẫu FT-05 ($n=0, n=12$), dấu, số mũ cơ sở và hai cờ cắt. &
Quét đủ 13 mức $n$ trên 1.634.048 giao dịch trước Packer, khớp oracle 100\%. \\
Nhận và giữ ngữ cảnh & Hồi quy hỗn hợp hai lõi trên toàn bộ 26 tổ hợp \texttt{OPS/n}. &
100.000 hoàn tất + 10.000 xác nhận sau tách S1 mỗi lõi, $0$ sai lệch. \\
Hỗ trợ backpressure & Hồi quy \textit{stall} ngẫu nhiên và trace giữ dữ liệu qua 3 cạnh liên tiếp. &
Đầu ra giữ nguyên trạng thái đến khi \texttt{out\_ready=1}, $0$ sai lệch. \\
Reset và phục hồi & Hủy giao dịch đang xử lý, không rò response cũ và nhận lại bình thường. &
Toàn bộ các ca hủy đều sạch dữ liệu cũ và phục hồi đúng ở giao dịch sau. \\
MAC ghép cặp NF/F & So \texttt{d/flags} trên cùng bộ $(A,B,C)$ với mô hình từng lõi. &
Hoàn tất 377.032 bộ dữ liệu tại \texttt{OPS1/n12}, $0$ sai lệch. \\
Tích lũy chuỗi dài & Đối chiếu từng bước NF/F, thanh ghi $Q$/count và lệnh \texttt{CLEAR/READ}. &
Hoàn tất 48 chuỗi (16.320 MAC/lõi + 96 lệnh điều khiển quire), $0$ sai lệch. \\
\bottomrule
\end{tabularx}
\end{engineeringtable}

\section{Tiêu chí quản lý độ chính xác và tính nhất quán kiểm chứng}

Để bảo đảm giá trị kiểm chứng không bị pha trộn giữa các cấu hình, quá trình
chạy hồi quy tuân thủ quy trình kiểm soát phiên bản chặt chẽ từ tệp RTL đầu vào,
mô hình tham chiếu, bộ sinh vector tới báo cáo đầu ra. Mỗi khi thành phần tổ
hợp được tinh chỉnh (như bước tách tải fanout 4 nhánh ở tầng S1), bộ kiểm tra
tương đương đơn vị (201.764 mẫu bit-exact cho khối \texttt{posit\_add\_prepare\_comb})
và bộ hồi quy ModelSim toàn cục đều được tái chạy trên cùng bộ dữ liệu hạt giống
cố định. Nhờ vậy, hệ thống kiểm chứng xác nhận thiết kế sau tối ưu bảo toàn tính
tương đương từng chu kỳ (\textit{cycle-exact}) so với RTL tham chiếu.

\section{Thiết kế tập kiểm và phân tầng kích thích}

Lượt hồi quy chức năng kết hợp xen kẽ các giao dịch độc lập và các chuỗi hồi tiếp.
Do lấy mẫu ngẫu nhiên thuần túy theo từ mã ít chạm tới các chuỗi regime rất dài
hoặc fraction có trọng số bit cao, bộ sinh kích thích thực hiện phân tầng có
chủ đích theo độ dài regime của hai toán hạng $(A,B)$ và bốn tổ hợp dấu $(s_A,s_B)$,
tạo thành lưới $8\times 8\times 4 = 256$ ô kiểm chứng cân bằng. Bên cạnh đó, các
cặp mẫu ép cận triệt tiêu (\textit{near-cancellation}) được chèn bổ sung để kích
thích mạnh các đường dịch chuẩn hóa dài tại bộ cộng \texttt{S1--S5} của lõi \texttt{MAC-F}.

\begin{engineeringtable}{Phân loại các biến đếm giám sát trong môi trường kiểm chứng.}{tab:counter-meaning}
\begin{tabularx}{\linewidth}{@{}>{\raggedright\arraybackslash}p{36mm}Y Y@{}}
\toprule
Loại biến đếm & Điều kiện ghi nhận & Vai trò trong đánh giá \\
\midrule
Accepted & \texttt{in\_valid} và \texttt{in\_ready} cùng bằng 1 tại cạnh lên. & Đối soát cân bằng giao thức ($\mathrm{Accepted} = \mathrm{Completed} + \mathrm{Aborted}$). \\
Completed & Giao dịch hoàn tất \textit{retire} và khớp dữ liệu/cờ với oracle. & Đánh giá khối lượng kiểm chứng hợp lệ. \\
Reset-aborted & Giao dịch đang chạy bị ngắt bởi tín hiệu \texttt{rst\_n=0}. & Kiểm chứng cơ chế hủy sạch và chống rò rỉ dữ liệu cũ. \\
Functional bin & Giao dịch hợp lệ kích hoạt đúng nhóm cấu hình/dữ liệu. & Bảo đảm 100\% không gian trạng thái chức năng được kiểm thử. \\
Structural point & Điểm đếm dòng lệnh, nhánh rẽ và chuyển trạng thái bit. & Đánh giá mức độ bao phủ phần cứng của RTL. \\
\bottomrule
\end{tabularx}
\end{engineeringtable}

Ngoài kiểm chứng toàn cục ở biên ngoài IP, quy trình kiểm thử phân lớp còn kiểm
tra trực tiếp tại các ranh giới trung gian gồm biên giải mã Parser (xử lý
\texttt{regime} cực trị và \texttt{Zero}/\texttt{NaR}), biên ngõ ra bộ nhân
(xác nhận giá trị Q24 và cờ cắt trước khi đi vào bộ cộng) và biên Packer cuối
(kiểm tra các ca hòa RNE, bão hòa biên và mã hóa số âm).
